\paragraph{Leading angular feedback on core-charge transport.}
A separate second-variation calculation follows
$\phi_\varepsilon=B+\varepsilon\eta Y_{20}+\varepsilon^2\Xi+O(\varepsilon^3)$
and evolves the forced spherical mean $\zeta=\langle\Xi\rangle_\Omega$.
Here $\Xi$ is half the second derivative, and $\zeta$ is an angular mean,
not the coefficient of normalized $Y_{00}$. Only this mean is needed
alongside the first tangent for the angle-integrated charge through
second order. The initial family is the volume-preserving ellipsoidal
deformation $r_\varepsilon^2=e^{-2c\varepsilon}(x^2+y^2)+e^{4c\varepsilon}z^2$,
$c=\sqrt{5/(16\pi)}$, with an explicitly differentiated finite-grid
normalization fixing total charge. For $Q_2$, the coefficient of
$\varepsilon^2$, define
\[
 b=\frac{Q_{2,\mathrm{core}}(32)-Q_{2,\mathrm{core}}(0)}{Q_0}
\]
in the same fixed core $r<6$, box $R=80$ and time window.
The three $(h,\Delta t)=(0.1,0.02),(0.05,0.02),(0.1,0.01)$ calculations
give $b=0.08737135,\ 0.08736583,\ 0.08737135$, respectively;
the spatially refined value exceeds the predeclared heuristic guard
$2.864\times10^{-5}$. This guard combines measured spatial, temporal and
charge-balance sensitivities with a fixed floor; it is not a certified
error bound or confidence interval. Stage-integrated physical charge flux
checks the core-charge difference, rather than defining it from the endpoint.

The interpretation requires the initial redistribution
(Fig.~\ref{fig:angular-second-order}). On the finer spatial grid,
$Q_{2,\mathrm{core}}(0)/Q_0=-0.10587389$, whereas
$Q_{2,\mathrm{core}}(32)/Q_0=-0.01850805$.
The radial reference itself has positive net core inflow over this window:
its core fraction increases from $0.58949606$ to $0.65267983$.
The positive second-order net-transport coefficient recovers much of an
initial core-charge deficit but leaves the final correction negative;
it does not show more final charge retained than in the radial reference.
The sampled response is nonmonotonic, with a transient positive core
correction before the negative final value; the decision uses the fixed
endpoint $T=32$, not a selected intermediate maximum.
The preparation fixes charge, not energy:
$E(\varepsilon,0)=E_0+\varepsilon^2 E_2+O(\varepsilon^3)$ with
$E_2/E_0=0.0224455$ on that grid. These are second-order Taylor coefficients at zero deformation,
not a finite-amplitude nonlinear three-dimensional evolution, a binding
comparison at equal energy, or evidence of stationary Q-ball formation.
As a subsequent analytical interpretation, the untruncated continuum
family has $E_{\mathrm{grad}}(\varepsilon)=
E_{\mathrm{grad},0}[2e^{-2c\varepsilon}+e^{4c\varepsilon}]/3$,
while its kinetic and potential integrals remain unchanged. Thus its
positive $E_2=4c^2E_{\mathrm{grad},0}$ is a geometric preparation cost;
this continuum identity neither replaces the stored discrete $E_2$ nor
constitutes a new acceptance criterion or determines the transport sign.

\begin{figure}[htbp]
\centering
\includegraphics[width=.98\linewidth]{angular-second-order-transport.pdf}
\caption{Second-order angular response from stored samples. Left:
$Q_{2,\mathrm{core}}(t)/Q_0$ for all three discretizations, with the initial
fine-space correction marked separately. Right: the initial-subtracted
coefficient $b(t)$ and the negative, independently stage-integrated
outward charge flux divided by $Q_0$, on the fine-space grid.
These routes check the same charge balance, not separate
physical replications. Straight lines join 41 stored samples; no finite
$\varepsilon$ is reconstructed. The final core correction remains negative
despite positive net transport. The endpoint guard is heuristic, not a
confidence interval, certified error or time-uniform uncertainty band.}
\label{fig:angular-second-order}
\end{figure}

% Numerical result independently reviewed: RESULT-REVIEW.txt cee669eda28b331f143750ed01913dc63e985b13ace98badac8e0c3458a652d3.
\paragraph{Compensating the initial energy through second order.}
A further, separately specified family keeps the same first angular
variation but replaces the initial chirp by
$\beta(\varepsilon)=\beta_0+\beta_2\varepsilon^2$.
On each grid, $\beta_2=-E_2/(\partial_\beta E_0)$ is determined from the
preparation energy, without using the later transport. The normalized
family fixes total charge exactly and cancels the initial energy change
through order $\varepsilon^2$; it does not assert exactly equal energies
at finite deformation. Only the homogeneous radial chirp tangent is newly
evolved, and its contribution is combined with the stored angular
second variation. The earlier fixed-charge result remains unchanged.

On the spatially refined grid, $\beta_2=-0.0082897517$ and the normalized
net-transport derivative with respect to $\beta$ is $-11.0018872$.
Their product adds $0.0912029133$ to the previous transport coefficient:
$b_{QE}=0.1785687477$, compared with $b_Q=0.0873658344$.
The base and time-refined values are $0.1785384553$ and $0.1785384478$;
the predeclared heuristic guard for the spatially refined transport
coefficient is $1.525216\times10^{-4}$, not a certified error interval.
The initial core correction is unchanged, while the final core coefficient
$c_{QE}=Q_{2,\mathrm{core},QE}(32)/Q_0$ is now $+0.0726948606$,
compared with $-0.0185080526$ for the previous family.
This final coefficient is reported separately from the
net-transport decision; no additional endpoint-success gate was introduced.
The compensation also reduces the initial inward radial current through
its negative chirp correction. Relative to the radial reference, this
energy-compensated family changes geometry and radial phase together;
its transport response cannot isolate geometry alone.
These coefficients concern the fixed endpoint $T=32$ in the same finite
box, not a finite-amplitude deformed solution, stationary formation or binding.
Figure~\ref{fig:fixed-energy-transport} separates the core coefficients
from their initial-subtracted transport for the two preparations.

\begin{figure}[htbp]
\centering
\includegraphics[width=.98\linewidth]{fixed-energy-transport.pdf}
\caption{Stored spatially refined second-variation coefficients with and
without initial-energy compensation. Left: the core-charge correction;
right: its change from the common initial correction. Straight segments
join 41 stored samples, without fitting or a new evolution. The original
family fixes charge; the compensated family also fixes initial energy
through order $\varepsilon^2$ while changing the radial chirp. Neither
curve represents a finite deformation or a percentage increase. Endpoint
caps on the transport panel are heuristic guards, not confidence intervals,
core-coefficient error bars or time-uniform bounds. The sign change of the
final core coefficient is not a binding or stability result.}
\label{fig:fixed-energy-transport}
\end{figure}

% Stored-state local-clock diagnostic, independently reviewed; finite scope only.
\paragraph{A local two-observable time-shift test.}
The energy-compensated response need not describe a new stationary state:
even an increased endpoint core coefficient could reflect a shifted phase
of the transient concentration. A fixed stored-state check at $T=32$
therefore compares two radial observables,
$C=Q_{\mathrm{core}}/Q_0$ and $H=R_{\mathrm{RMS}}^2/6^2$, where $R_{\mathrm{RMS}}^2$
is the second radial moment weighted by $|\phi|^2$ and divided by its
spatial integral. Write their second-order coefficients as $C_2,H_2$.
Here $C_2(32)=c_{QE}$, not the initial-subtracted $b_{QE}$.
A common local shift $t\mapsto t+\varepsilon^2\tau$ would require
$(C_2,H_2)=\tau(\dot C_0,\dot H_0)$, hence
\[
 \mathcal D=\dot C_0 H_2-\dot H_0 C_2=0.
\]
The moment and its derivative were evaluated directly from the stored
Cauchy data, including the norm-denominator variation; no time-series fit
or additional field evolution was used. On the fine spatial grid,
$H_2=-1.0936132$ and $\mathcal D=0.0021220476$, exceeding the predeclared
heuristic sensitivity guard $6.820019\times10^{-5}$ in inverse model-time
units. Thus this pair is inconsistent with one common local clock shift
under that diagnostic rule. The charge-only shift coefficient is
$\tau_C=C_2/\dot C_0=-9.1439597$, after its separate denominator check:
it multiplies $\varepsilon^2$ and is not a measured finite time shift.
This is only a two-observable endpoint comparison, conditional on the
existing variational outputs; it establishes neither persistence of the
core enhancement nor binding, stationary formation or a full-field distance.
The sensitivity guard is not a certified error bound.
